% Chapter 6, typeset from lectures/L06: the README is the opener and the closing review, the
% lecture's Appendix A is a section with the same subsections, and its Appendix B, the exercises,
% closes the chapter. The serial_driver_demo is typeset straight from its files in §6.1.6.

\chapter{Serial Communication (UART)}
\label{c6:ch}

{\large\itshape A real channel to the PC, one bit at a time.\par}

\begin{chapterbox}{In this chapter}
\begin{itemize}
  \item Why serial: a real communication channel to a PC instead of blink-and-infer debugging.
  \item Configuring UART: \code{UCSR0B}, \code{UCSR0C}, and the \code{UBRR0} baud-rate formula.
  \item Sending data: \code{UDR0}, the \code{UDRE0} flag, and building a reusable text-transmit
    helper.
  \item Mixing text and numbers with \code{sprintf()}.
  \item Receiving data: \code{RXEN0}, the \code{RXC0} flag, and reading \code{UDR0} instead of
    writing it.
  \item Worked example: button toggles an LED and reports its new state over serial, using a
    driver split across its own header and source file.
\end{itemize}
\end{chapterbox}

This chapter gives the board a way to talk back. Until now every program has reported what it was
doing by blinking an LED; with the UART it can send text to a PC. After the chapter you should be
able to explain what baud rate is and compute \code{UBRR0} for a given clock and baud rate, write a
byte and then a whole string over the UART, combine \code{sprintf()} with transmission to report a
mix of text and numbers, explain why the code waits on \code{UDRE0} before writing a new byte, and
receive a byte and explain how \code{RXC0} mirrors \code{UDRE0}.

Read \secref{c6:sec:uart}, the substance of the chapter, in order: transmit first, then the
chapter's demo, the book's first program split across several files, then receiving. Then work
the exercises in \secref{c6:sec:exercises}. The Reflection and Calculation exercises are answered
in Appendix~\ref{app:answers}; the Code exercises have no printed solution.

% Section 6.1, from lectures/L06/appendix/a_uart.md.

\section{Serial Communication (UART)}
\label{c6:sec:uart}\label{c6:app:a}

\subsection{Why serial}
\label{c6:a:1}
\begin{itemize}
  \item Every debugging technique so far has been ``blink an LED and infer what happened.'' UART
    (Universal Asynchronous Receiver/Transmitter) gives you a real communication channel to a PC,
    \code{printf()}-style debugging, logging sensor data, or eventually two-way communication with
    a host application.
  \item It's called \emph{serial} because it sends one bit at a time over a single wire, framed by
    a start bit and one or more stop bits per byte, cheap on wiring, at the cost of throughput.
    This chapter builds up transmit first (\secref{c6:a:2}--\ref{c6:a:6}), then receiving
    (\secref{c6:a:7}), which reuses almost the same registers.
  \item Speed is measured in \textbf{baud rate}, bits per second. 9600 baud is the traditional
    default and what we'll use here.
\end{itemize}

\subsection{Configuring UART}
\label{c6:a:2}
Three registers:
\begin{enumerate}
  \item \textbf{\code{UCSR0B}} enables transmit (and, later, receive):

\begin{ccode}
UCSR0B = (1U << TXEN0);   // Enable transmission.
\end{ccode}

  \item \textbf{\code{UCSR0C}} sets the frame format. 8 data bits, no parity, 1 stop bit (the
    near-universal default) needs:

\begin{ccode}
UCSR0C = (1U << UCSZ00) | (1U << UCSZ01);
\end{ccode}

  \item \textbf{\code{UBRR0}} holds the baud rate, computed from the CPU clock:

\begin{console}
UBRR0 = F_CPU / (16 * baud_rate) - 1
\end{console}

    \noindent At 16\,MHz for 9600 baud, that's \code{103}:

\begin{ccode}
UBRR0 = 103U;
\end{ccode}
\end{enumerate}

\begin{ccode}
#define BAUD 103U // UBRR value for 9600 bps with a 16 MHz clock.

/**
 * @brief Initialize UART with a 9600 bps baud rate.
 */
static void serial_init(void)
{
    // Enable UART transmission.
    UCSR0B = (1U << TXEN0);

    // Set character size to eight bits.
    UCSR0C = (1U << UCSZ00) | (1U << UCSZ01);

    // Set baud rate to 9600 bps.
    UBRR0 = BAUD;
}
\end{ccode}

\subsection{Sending data}
\label{c6:a:3}
Transmission happens one byte at a time through the \code{UDR0} data register. Before writing a
new byte, wait for the previous one to move out of \code{UDR0} (into the transmit shift register,
which sends it bit by bit); signalled by the \code{UDRE0} (Data Register Empty) flag in
\code{UCSR0A}:

\begin{ccode}
/**
 * @brief Transmit a single byte over UART, blocking until the data register is free.
 *
 * @param[in] byte Byte to transmit.
 */
static void serial_write_byte(const uint8_t byte)
{
    while (0U == (UCSR0A & (1U << UDRE0))) {} // Wait for the register to free up.
    UDR0 = byte;
}
\end{ccode}

\noindent Sending a whole string is just calling that in a loop, plus a trailing \code{\r\n}
(carriage return, then line feed) so each message starts on a fresh line in the terminal:

\begin{ccode}
/**
 * @brief Transmit text over UART.
 *
 * @param[in] text Text to transmit.
 */
static void serial_print(const char* text)
{
    // Terminate if the text pointer is NULL.
    if (NULL == text) { return; }

    // Transmit each character one by one.
    for (uint16_t i = 0U; text[i]; ++i)
    {
        serial_write_byte((uint8_t)(text[i]));
    }
    serial_write_byte('\r');
    serial_write_byte('\n');
}
\end{ccode}

\subsection{Simple TX example}
\label{c6:a:4}
A first fully working example: \secref{c6:a:2}'s \code{serial_init()} and \secref{c6:a:3}'s
\code{serial_write_byte()}/\allowbreak\code{serial_print()} combined into a single flashable program. LED1 on
D9 blinks every 500\,ms, and each transition is reported over serial. \code{delay_ms()} is the
same hand-rolled helper built on \code{_delay_ms()} from \secref{c2:a:6}, reused rather than
redefined from scratch. Everything still lives in one file here; \secref{c6:a:6} splits the driver
out into its own header and source file once there's enough of it to be worth the split.

\begin{ccode}
#define F_CPU 16000000UL // CPU frequency in Hz.

#include <stddef.h>
#include <stdint.h>

#include <avr/io.h>
#include <util/delay.h>

/** GPIO pins. */
#define LED1 1U // D9  -> PORTB1.

/** GPIO operations. */
#define LED1_TOGGLE PINB = (1U << LED1)    // Toggle LED1.
#define LED1_ENABLED (PORTB & (1U << LED1)) // High if LED1 is enabled, low otherwise.

/** Time parameters. */
#define BLINK_SPEED_MS 500U // Blink speed in ms.
#define TICK_1MS 1U         // Generate a 1 ms tick in delay_ms().

/** Serial parameters. */
#define BAUD 103U            // UBRR value for 9600 bps with a 16 MHz clock.
#define NEW_LINE '\n'        // New line character.
#define CARRIAGE_RETURN '\r' // Carriage return.

/**
 * @brief Initialize UART with a 9600 bps baud rate.
 */
static void serial_init(void)
{
    // Enable UART transmission.
    UCSR0B = (1U << TXEN0);

    // Set character size to eight bits.
    UCSR0C = (1U << UCSZ00) | (1U << UCSZ01);

    // Set baud rate to 9600 bps.
    UBRR0 = BAUD;
}

/**
 * @brief Transmit a single byte over UART, blocking until the data register is free.
 *
 * @param[in] byte Byte to transmit.
 */
static void serial_write_byte(const uint8_t byte)
{
    // Wait until the transmit buffer is ready for a new byte.
    while (0U == (UCSR0A & (1U << UDRE0))) {}

    // Place the new byte in the transmit data register.
    UDR0 = byte;
}

/**
 * @brief Transmit text over UART.
 *
 * @param[in] text Text to transmit.
 */
static void serial_print(const char* text)
{
    // Terminate if the text pointer is NULL.
    if (NULL == text) { return; }

    // Transmit each character one by one.
    for (uint16_t i = 0U; text[i]; ++i)
    {
        serial_write_byte((uint8_t)(text[i]));
    }
    serial_write_byte('\r');
    serial_write_byte('\n');
}

/**
 * @brief Generate delay.
 *
 * @param[in] ms Delay duration in ms.
 */
static void delay_ms(const uint16_t ms)
{
    for (uint16_t i = 0U; i < ms; ++i)
    {
        _delay_ms(TICK_1MS);
    }
}

/**
 * @brief Set up system.
 */
static void setup(void)
{
    // Configure LED1 as output.
    DDRB = (1U << LED1);

    // Initialize serial driver.
    serial_init();
}

/**
 * @brief Application entry point.
 *
 * @return 0 on termination of the program (should never occur).
 */
int main(void)
{
    setup();

    // Toggle LED1 and print its new state every 500 ms.
    while (1)
    {
        LED1_TOGGLE;
        if (LED1_ENABLED) { serial_print("LED1 enabled!"); }
        else { serial_print("LED1 disabled!"); }
        delay_ms(BLINK_SPEED_MS);
    }
    return 0;
}
\end{ccode}

\noindent\code{LED1_ENABLED} reads \code{PORTB} here, not \code{PINB}. A pin's level reaches
\code{PINB} through a synchronizer that runs one clock cycle behind, so reading back a level the
program has just set needs an instruction in between (the datasheet uses a \code{nop}). In this
loop the compiler puts the read directly after the toggle, and a \code{PINB} read would still see
the old level, reporting every transition the wrong way round. \code{PORTB} holds what the program
last wrote, and changes at once.

\subsection{Mixing text and numbers}
\label{c6:a:5}
\code{sprintf()} (from \code{<stdio.h>}) formats into a buffer exactly like \code{printf()} formats
to the terminal; build the string there, then hand it to \code{serial_print()}. A full program,
extending \secref{c6:a:4} with a counter reported once every 500\,ms instead of a fixed message:

\begin{ccode}
#define F_CPU 16000000UL // CPU frequency in Hz.

#include <stddef.h>
#include <stdint.h>
#include <stdio.h>

#include <avr/io.h>
#include <util/delay.h>

/** Time parameters. */
#define REPORT_INTERVAL_MS 500U // Report interval in ms.
#define TICK_1MS 1U             // Generate a 1 ms tick in delay_ms().

/** Serial parameters. */
#define BAUD 103U // UBRR value for 9600 bps with a 16 MHz clock.

/** Message parameters. */
#define MSG_LEN 50U // Max length of a formatted message, including the NUL terminator.

/**
 * @brief Initialize UART with a 9600 bps baud rate.
 */
static void serial_init(void)
{
    // Enable UART transmission.
    UCSR0B = (1U << TXEN0);

    // Set character size to eight bits.
    UCSR0C = (1U << UCSZ00) | (1U << UCSZ01);

    // Set baud rate to 9600 bps.
    UBRR0 = BAUD;
}

/**
 * @brief Transmit a single byte over UART, blocking until the data register is free.
 *
 * @param[in] byte Byte to transmit.
 */
static inline void serial_write_byte(const uint8_t byte)
{
    // Wait until the transmit buffer is ready for a new byte.
    while (0U == (UCSR0A & (1U << UDRE0))) {}

    // Place the new byte in the transmit data register.
    UDR0 = byte;
}

/**
 * @brief Transmit text over UART.
 *
 * @param[in] text Text to transmit.
 */
static void serial_print(const char* text)
{
    // Terminate if the text pointer is NULL.
    if (NULL == text) { return; }

    // Transmit each character one by one.
    for (uint16_t i = 0U; text[i]; ++i)
    {
        serial_write_byte((uint8_t)(text[i]));
    }
    serial_write_byte('\r');
    serial_write_byte('\n');
}

/**
 * @brief Generate delay.
 *
 * @param[in] ms Delay duration in ms.
 */
static void delay_ms(const uint16_t ms)
{
    for (uint16_t i = 0U; i < ms; ++i)
    {
        _delay_ms(TICK_1MS);
    }
}

/**
 * @brief Set up system.
 */
static void setup(void)
{
    // Initialize serial driver.
    serial_init();
}

/**
 * @brief Application entry point.
 *
 * @return 0 on termination of the program (should never occur).
 */
int main(void)
{
    uint16_t counter = 0U;
    setup();

    // Report an incrementing counter every 500 ms.
    while (1)
    {
        // Parse the counter value into a string:
        char msg[MSG_LEN] = {'\0'};
        sprintf(msg, "Count: %u", counter++);

        // Transmit the string content over UART, then wait 500 ms.
        serial_print(msg);
        delay_ms(REPORT_INTERVAL_MS);
    }
    return 0;
}
\end{ccode}

\noindent\code{\%d} for \code{int}, \code{\%u} for unsigned, \code{\%u} used here since
\code{counter} is a \code{uint16_t}, which is an \code{unsigned int} on this target.
Floating-point formatting in AVR's default \code{sprintf()} is often disabled or costly to link
in; prefer integers (e.g.\ scale a reading to hundredths and print the two parts separately)
unless you've confirmed float formatting is available in your toolchain.

\subsection{Example: button toggles an LED, reports over serial}
\label{c6:a:6}
Button on D13, LED on D9. Each press flips the LED and reports its new state.

\bookimage{lectures/L06/appendix/images/circuit_serial_with_gpios.png}{The serial demo's
circuit: a button on D13 and an LED on D9, with the board's USB port as the serial link to the
PC.}{c6:fig:serial}

This is also the book's first multi-file example: rather than inlining every function here, the
serial driver lives in its own header and source file, and \code{main.c} \code{#include}s it like
any other driver:

\begin{console}
serial_driver_demo/
├── Makefile
├── main.c
├── include/
│   └── driver/
│       └── serial.h
└── source/
    └── driver/
        └── serial.c
\end{console}

\noindent The driver (\file{driver/serial.h} and \file{driver/serial.c}, both printed in full
below) adds one function beyond \secref{c6:a:3}--\ref{c6:a:4}'s \code{serial_write_byte()}/%
\code{serial_print()}: \code{serial_write()}, which takes a raw \code{uint8_t} buffer and a
length, for arbitrary binary data rather than NUL-terminated text. \code{main.c} below still only
needs \code{serial_print()}, the same convenience wrapper already familiar from
\secref{c6:a:3}--\ref{c6:a:4}, with one difference: the driver's version doesn't append a line
ending of its own. It sends the text as written, turning each \code{\n} in it into \code{\r\n} on
the way out, which is why the messages below end in \code{\n}. \code{serial_write_byte()} is a
\code{static} helper inside \file{serial.c}, no longer something callers can reach. \code{main.c}
itself is shown in full after the driver.

\subsubsection{The driver in full}
\cfile{lectures/L06/serial_driver_demo/include/driver/serial.h}
\needspace{8\baselineskip}
\cfile{lectures/L06/serial_driver_demo/source/driver/serial.c}

\subsubsection{The program}
\cfile{lectures/L06/serial_driver_demo/main.c}

\noindent\code{ISR(PCINT0_vect)} only toggles LED1 and sets \code{led1_event}; the UART report
happens in \code{main()}'s loop instead, once per event, never inside the ISR itself.
\code{serial_print()} blocks byte-by-byte on \code{UDRE0} until each byte is sent, and an ISR that
spends that time transmitting holds up every other interrupt for the whole transmission, including
the debounce timer this very feature depends on. Keep ISRs to the minimum work needed (flip a pin,
set a flag) and let \code{main()}'s loop handle anything that can block; \secref{c7:a:4} reuses
this exact pattern for \code{eeprom_write()}, another slow, blocking call that doesn't belong in an
ISR either.

\subsubsection{Building and running}
\label{c6:sec:build}
On Linux, \code{cd} into \file{lectures/L06/serial_driver_demo} and run \code{make} to build, then
\code{make flash} to program the board, the same \file{Makefile} pattern as \secref{c2:d:8}:

\makefile{lectures/L06/serial_driver_demo/Makefile}

\noindent On Windows via Microchip Studio, there's no Makefile support: recreate this same folder
structure inside the project, then add \file{include} as an include directory under the project's
toolchain settings (Project Properties $\rightarrow$ Toolchain $\rightarrow$ AVR/GNU C Compiler
$\rightarrow$ Directories), the IDE's equivalent of the Makefile's \code{-Iinclude}.

Open a serial terminal (the Arduino IDE's Serial Monitor, or \code{screen}/\code{minicom} from the
command line) at 9600 baud on the Uno/Nano's port to see the output.

\subsection{Receiving data}
\label{c6:a:7}
Receiving is transmit's mirror image: one more enable bit in \code{UCSR0B}, one more status flag
in \code{UCSR0A} to poll, and the same \code{UDR0} data register, just read instead of written.
\begin{enumerate}
  \item \textbf{\code{RXEN0}} (Receiver Enable) sits right next to \code{TXEN0} in \code{UCSR0B},
    the register \secref{c6:a:2} introduced. Set both together to enable transmit and receive at
    once:

\begin{ccode}
UCSR0B = (1U << TXEN0) | (1U << RXEN0);   // Enable transmission and reception.
\end{ccode}

  \item \textbf{\code{RXC0}} (RX Complete), in \code{UCSR0A}, is \code{UDRE0}'s counterpart: it's
    clear while nothing new has arrived, and sets itself the moment a fully-received byte is
    sitting in \code{UDR0}, waiting to be read. Reading \code{UDR0} clears it automatically,
    exactly as reading \code{UDR0} also readies it for the next incoming byte.
  \item \textbf{\code{UDR0}} is the very same register \code{serial_write_byte()} writes to in
    \secref{c6:a:3}; reading it instead of writing it returns whatever byte just arrived.
\end{enumerate}
Put together, the receive-side equivalent of \code{serial_write_byte()} is:

\begin{ccode}
/**
 * @brief Receive a single byte over UART, blocking until one arrives.
 *
 * @return Byte received.
 */
static char serial_read_byte(void)
{
    while (0U == (UCSR0A & (1U << RXC0))) {} // Wait for a byte to arrive.
    return UDR0;
}
\end{ccode}

\noindent This blocks indefinitely if nothing ever arrives, fine for a first pass, but worth
keeping in mind once a program also needs to do other work while waiting. Reading more than one
byte into a buffer, with an optional timeout so a silent sender can't hang the program forever,
builds directly on \code{serial_read_byte()} the same way \secref{c6:a:6}'s driver
\code{serial_write()} builds on \code{serial_write_byte()}.

\code{read()} below does exactly that: it fills a caller-supplied buffer up to \code{buflen} bytes,
and takes an optional \code{timeout_ms} so a sender that never shows up can't hang the program
forever.

\begin{ccode}
/**
 * @brief Receive up to buflen bytes over UART into buf.
 *
 * @param[out] buf Buffer to receive into.
 * @param[in] buflen Size of buf in bytes.
 * @param[in] timeout_ms Milliseconds to wait for buflen bytes before giving up.
 *            0 blocks indefinitely.
 *
 * @return Number of bytes received, or -1 if buf is NULL or buflen is 0.
 */
int16_t read(uint8_t* buf, const uint16_t buflen, const uint16_t timeout_ms)
{
    // Check the input parameters, return -1 if invalid.
    if ((NULL == buf) || (0U == buflen)) { return -1; }
    uint16_t rx_count = 0U;

    if (0U == timeout_ms)
    {
        // Read indefinitely until the buffer is full if no timeout has been specified.
        while (buflen > rx_count)
        {
            // Wait until a byte has been received.
            while (0U == (UCSR0A & (1U << RXC0))) {}
            buf[rx_count++] = UDR0;
        }
    }
    else
    {
        for (uint16_t i = 0U; i < timeout_ms; ++i)
        {
            while ((buflen > rx_count) && (UCSR0A & (1U << RXC0)))
            {
                buf[rx_count++] = UDR0;
            }
            // Stop reading if the read buffer is full.
            if (buflen <= rx_count) { break; }

            // Wait a millisecond before reading again.
            delay_ms(1U);
        }
    }
    // Return the number of bytes read.
    return (int16_t)(rx_count);
}
\end{ccode}

\noindent\code{read()} extends \code{serial_read_byte()} from a single byte to a whole buffer,
with \code{timeout_ms} choosing between two modes:
\begin{itemize}
  \item \textbf{\code{timeout_ms == 0}}: block indefinitely, the same as
    \code{serial_read_byte()}, until \code{buflen} bytes have arrived.
  \item \textbf{\code{timeout_ms} nonzero}: give up after roughly that many milliseconds even if
    the buffer isn't full yet, pacing itself with the \code{delay_ms()} helper first defined in
    \secref{c2:a:6}, one millisecond at a time, rather than redefining it here.
\end{itemize}
\code{buf} is checked against \code{NULL} before use, the same reasoning as \code{vector_resize()}
back in \secref{c1:a:11}: this is written as a reusable, exported function with call sites outside
this file, not a private helper like the \code{delay_ms()} of \secref{c3:a:4}, so the check is
doing real work rather than just demonstrating the habit.

The return value distinguishes three outcomes: \code{-1} for a bad call (\code{NULL} buffer or
\code{buflen == 0}), the number of bytes actually received if the timeout elapsed first, or
\code{buflen} itself once the buffer filled completely.

% Section 6.2, from lectures/L06/README.md: the goals, the evaluation questions, and the next
% lecture.

\needspace{12\baselineskip}
\section{Review}
\label{c6:sec:review}

\needspace{6\baselineskip}
\subsection*{What you should be able to do now}
After this chapter, you should be able to:
\begin{itemize}
  \item Explain what baud rate is and compute \code{UBRR0} for a given clock and baud rate.
  \item Write a byte, and then a whole string, over UART.
  \item Combine \code{sprintf()} with UART transmission to report a mix of text and numeric data.
  \item Understand why the code waits on \code{UDRE0} before writing a new byte.
  \item Receive a byte over UART and explain how \code{RXC0} mirrors \code{UDRE0}.
\end{itemize}

\needspace{6\baselineskip}
\subsection*{Questions to test yourself}
\begin{enumerate}
  \item What does the \code{UDRE0} flag represent, and what would happen if
    \code{serial_write_byte()} didn't wait for it?
  \item Compute \code{UBRR0} for 19200 baud at a 16\,MHz clock.
  \item Why is \code{sprintf()} used to build a message before transmitting it, rather than trying
    to transmit an \code{int} directly?
  \item What does \code{RXC0} represent, and what would \code{serial_read_byte()} do if it didn't
    wait for it?
\end{enumerate}

\needspace{6\baselineskip}
\subsection*{Looking ahead}
This chapter closes the introductory arc:
\begin{itemize}
  \item C fundamentals (Chapter~\ref{c1:ch}).
  \item Digital I/O and bit manipulation (Chapter~\ref{c2:ch}).
  \item Reacting to events instead of polling for them (Chapter~\ref{c3:ch}).
  \item Reading the analog world (Chapter~\ref{c4:ch}).
  \item Timing without blocking (Chapter~\ref{c5:ch}).
  \item Reporting back to a PC (this chapter).
\end{itemize}
Chapter~\ref{c7:ch} rounds out the peripheral survey with EEPROM, and Chapter~\ref{c8:ch} covers
the watchdog timer. A follow-on arc on driver design starts after that: bundling a peripheral's
state into a struct instead of one-off macros and globals, building directly on the LEDs, buttons,
and timers from these six chapters.

% Section 6.3, from lectures/L06/appendix/b_exercises.md.

\section{Exercises}
\label{c6:sec:exercises}\label{c6:app:b}

\begin{aside}
\textbf{How to work these.} The Reflection and Calculation exercises are answered on paper, from
the chapter. The Code exercises are programs for the board, except Exercise~\ref{c6:ex:3}, which
is a plain PC program: build and flash the board programs the way \secref{c2:app:c} (Microchip
Studio) or \secref{c2:app:d} (Linux) describes, or run them in the simulator of
Appendix~\ref{sim:ch} if you have no board yet.

\textbf{How to check your work.} The Reflection and Calculation exercises are answered in
Appendix~\ref{app:answers}. The Code exercises have no printed solution: the board, or for
Exercise~\ref{c6:ex:3} the terminal, tells you whether a program does what the exercise describes.
Try each exercise before you read its answer.
\end{aside}

% ------------------------------------------------------------------------------------------
\exerciseset{Concepts}
\label{c6:set:1}

\calcexercise{\code{UBRR0} for 19200 baud}
\label{c6:ex:1}
Calculate the value of \code{UBRR0} for a baud rate of \code{19200} at a 16\,MHz clock. Round to
the nearest integer.

\reflexercise{What \code{UDRE0} means}
\label{c6:ex:2}
Explain what the \code{UDRE0} flag represents and why \code{serial_write_byte()} waits for it
before writing to \code{UDR0}.

% ------------------------------------------------------------------------------------------
\exerciseset{Program}
\label{c6:set:2}

\codeexercise{Decimal digits by hand}
\label{c6:ex:3}
Write a \code{serial_write_uint(const uint8_t byte)} helper that prints \code{byte} as decimal
text, one digit at a time, e.g.\ \code{serial_write_uint(243)} prints \code{243}. Extract each
digit yourself via division/modulo; don't just hand the whole number to
\code{printf("\%u", byte)}/\code{sprintf()} and let it do the conversion for you, that's the same
idea you'll need on any future toolchain where \code{sprintf()}/float formatting isn't available.

This is a console exercise: a plain PC program is enough, no board or UART required.
\code{printf()} (or \code{putchar()}) stands in for \code{serial_write_byte()}, one call per digit.

Hint: dividing/modulo-ing a number by \code{10} repeatedly gives you its digits least-significant
first; printing them in that order comes out backwards. Instead, find the largest power of ten
that still fits inside \code{byte} first (\code{uint8_t} maxes out at \code{255}, so at most 3
digits), then walk that divisor down by a factor of \code{10} each step, digits come out in the
right order with no need to reverse anything afterward. Watch out for \code{0}: make sure the
search for the starting divisor doesn't shrink it all the way down to nothing before the loop that
actually emits digits ever runs.

\codeexercise{Receiving a command}
\label{c6:ex:4}
Extend \file{serial_driver_demo}'s driver with receive support, following \secref{c6:a:7}'s
\code{RXEN0}/\code{RXC0} pattern: enable the receiver in \code{serial_init()}, then add a
\code{serial_read()} to \file{serial.h}/\file{serial.c} that blocks until a byte arrives.

Use it in \code{main()} so that receiving \code{'t'} over serial toggles LED1 the same way pressing
the button already does, reporting the new state the same way \code{main()} already reports a
button press.

\codeexercise{A temperature reporter}
\label{c6:ex:5}
Capstone, tying together ADC (Chapter~\ref{c4:ch}), timers (Chapter~\ref{c5:ch}), and UART (this
chapter). Read ambient temperature from a TMP36 sensor on A0 and report it over UART, either once a
minute (Timer-driven) or immediately on a button press on D13, debounced with a timer so contact
bounce doesn't spam the report:
\begin{itemize}
  \item Reuse \code{adc_read()} from \secref{c4:a:3} to sample A0.
  \item TMP36 outputs \code{750 mV} at \code{25 °C} and rises \code{10 mV} per \code{°C};
    converting the ADC's \code{0-1023} reading to a voltage and then to a temperature gives:

\begin{console}
T (°C) = 500 * adc_reading / 1023 - 50
\end{console}

  \item The reported message should look like \code{"Temperature: 23 C"} (build it with
    \code{sprintf()}, as in \secref{c6:a:5}).
\end{itemize}
For a harder version: keep the five most recent readings in a small array and report the running
average alongside the instantaneous reading; and instead of a fixed one-minute report interval,
track the time between the last five button presses and auto-adjust the report interval to match
how often the button is actually being pressed.

